% Shared by main.tex (article) and main_ieee.tex (IEEEtran).
Multi-agent architectures built from large language models are widely deployed,
but reports of their benefit conflict, and practitioners have no way to
anticipate whether one will help before building it. We argue two fixable
problems explain the disagreement: baselines are rarely matched on compute, and
no measurable task property has been identified that governs the outcome. We
show that a $k$-agent ensemble's accuracy decomposes exactly into
\emph{coverage}---the chance any of its candidates is correct---and
\emph{selection efficiency}, the fraction of that headroom its selector
converts. We then give a latent-score model that turns a cheap pairwise probe,
the \textbf{generator--verifier gap} $\gap$, into a parameter-free prediction of
selection efficiency. Across \nTasks{} tasks in \nFamilies{} families spanning
the verification spectrum, \nModels{} model scales and \nArch{} architectures,
with exact token accounting over \nGenerations{} generations, we find that
against a single agent priced at the same token budget, \nFlip{} of the apparent
multi-agent gains disappear and \fracMatchedNegative{} of all configurations are
net negative. Diagnostics measured on \nCalib{} calibration tasks predict held-out lift with
$R^2=\rTwoApriori{}$ (Pearson $r=\corrApriori{}$), with no fitted parameters and
without the multi-agent system ever being constructed. A case study on citation generation shows the same split in a deployed setting:
an external registry check attains selection efficiency $\citEtaAPI{}$ at zero
token cost, while the same model used as its own critic attains
$\citEtaLLM{}$. The pattern that survives is narrow: ensembles selected by an external program
verifier gain reliably, self-consistency gains where the answer distribution is
peaked on the truth, and ensembles selected by the model judging itself gain
nowhere. Multi-agent structure is not a general accuracy technique; it is a way
of spending compute that pays off exactly when checking an answer is easier
than producing one.
